% 从旧数学附录迁出的筛法与判素说明；不新增算法。
\section{杜教筛：预筛与缓存都不可省略}\label{knowledge-dujiao}\index{Du Jiao sieve}
整除反演建模（第~\ref{knowledge-mobius-model}~节，第~\pageref{knowledge-mobius-model}~页）的两条卷积恒等式给出 \texttt{DuJiao} 的递推，见代码中的
\texttt{mertens} 与 \texttt{totient\_sum}。
把商相同的一段合并，系数为区间长度，递归参数严格变小；0 直接查表返回 0。
超过预筛上限的结果在返回时写入缓存，后续相同参数直接复用。
WIDA 收集的 \texttt{sumMu} 示例查询了 \texttt{fMu}，但未写入新结果；
本库已有实现包含写回，不能照搬遗漏缓存的版本。

设最大查询为 $N$，预筛上限为 $B$。单次查询的常用上界为
$O(B+N/\sqrt B)$，其中哈希访问按期望常数处理；取 $B$ 约 $N^{2/3}$
得到 $O(N^{2/3})$。多个不同查询可能继续增加缓存，不能把单次空间估计当成任意多询问的上界。
保留缓存可复用，主动清空仅影响重算开销；预筛表与上限应视为只读。

接口范围为 $0\le n\le10^{11}$，返回精确整数；欧拉函数前缀和用 128 位。
WIDA 的 \texttt{Z} 模数结果可在精确莫比乌斯前缀和返回后取模得到。
负值先求余，仅在余数为负时加模数，避免无条件加模数的溢出风险。

两套实现以独立试除构建到 200000 的前缀参考，检查四种预筛上限、零输入、
乱序与重复查询、缓存写回及清空重建。
上限 $10^{11}$ 使用 $10^6$ 与 $2\times10^6$ 两种预筛规模核对双版本结果；
这一项属于边界运行与一致性证据，不是独立答案对拍。
普通与 ASan/UBSan 均通过，暂无该接口的在线 AC。

\section{64 位判素与质因数分解}\label{knowledge-miller-bounds}\index{Miller--Rabin}\index{Pollard--Rho}
本库判素采用七个固定底数的 Miller--Rabin，覆盖无符号 64 位整数：
\[
2,\ 325,\ 9375,\ 28178,\ 450775,\ 9780504,\ 1795265022.
\]
将 $n-1=d2^s$、$d$ 为奇数，检查 $a^d$ 及后续平方链。
小于 2 的输入不视为素数；底数模 $n$ 为 0 时跳过该底数，不可误判为合数。
模乘用无符号 128 位容纳完整的两个 64 位数之积。
固定底数的适用范围依据公开 SPRP 记录，不外推到任意精度整数。

\url{https://miller-rabin.appspot.com/}

Pollard--Rho 用 $f(x)=x^2+c\pmod n$ 迭代，并求 $\gcd(|x-y|,n)$。
只有 $1<d<n$ 的因子才返回；一次尝试超过轮数上限或得到 $d=n$，就换参数重试。
其中模平方后的加法也在 128 位中执行，避免 $n$ 接近 $2^{64}$ 时加 $c$ 溢出。
该过程保证返回真因子，但随机尝试次数没有确定的最坏上界。
递归分解两侧，直到判为素数，最后排序并保留重复因子。

调用 \texttt{factor(1)} 得到空数组；对 $n\ge2$，返回单元素说明 $n$ 为素数。
POJ 1811 所需的最小质因子可取排序结果首项；P4718 所需的最大质因子可取末项。
不要对空数组取首尾，也不要绕过 \texttt{factor} 把素数传给内部 \texttt{rho}。

双版本对百万范围以独立埃氏筛核验判素，对 $1\ldots10000$ 以试除核验完整因子列表。
另验证三个随机种子下的大质数平方、重复因子、强伪素数和 $2^{64}-1$；
期望因子的素性均以试除独立验证。
$2^{61}-1$ 使用完整分解 $n-1$ 的 Lucas 素性证书检查；模乘与模幂另以任意精度运算对拍。
完整 POJ 1811 与 P4718 驱动分别验证 350 组输入，检查素数标签与最小、最大质因子。
P4718 压力批次含 300 个接近 $10^{18}$ 的平衡半素数，本机普通编译运行约 1.99、1.83 秒；
机器差异与随机路径都会影响耗时，本地运行时间不直接代表线上性能。
两套代码随后通过 P4718 全部 20 个测试点，最长单点显示约 2.11、2.12 秒，低于当前页面的 4 秒限制。
该 AC 检查素数标签与最大质因子，范围为 $2\le n\le10^{18}$。
完整因子重数、最小质因子、$n=1$ 与更大的无符号 64 位输入仍依据本地验证；
普通与 ASan/UBSan 均通过，不构成随机运行时间保证。

来源：kuangbin 2.6；固定底数记录见上方链接。

\section{单个数的欧拉函数}\label{knowledge-phi-single}\index{Euler totient}
实现见第~\pageref{compact-euler_phi}~页。对 $n\ge1$，由容斥得到
\[
\varphi(n)=n\prod_{p\mid n}\left(1-\frac1p\right),
\]
其中只枚举不同质因子，故质数幂 $p^k$ 的贡献为 $p^{k-1}(p-1)$。
每次更新按先除后乘的顺序：$\mathrm{ans}\gets(\mathrm{ans}/p)(p-1)$。
当前尚未处理的质因子仍整除答案，而且答案只会减小，乘法结果不会超过原始 $n$。

无分解器参数的入口采用试除，循环上界用 $p\le n/p$，剩下的 $n>1$ 再处理一次。
这能避免 $p^2$ 溢出，但不消除 $O(\sqrt n)$ 的耗时；不适合试除 64 位大素数。
传入现有 \texttt{PollardRho} 对象时先分解，再在排序后的因子列表上去重更新；
不重复构造分解器，后续调用沿用其随机状态。该入口没有确定的最坏运行时间。
$\varphi(1)=1$，两个入口均接受完整的正无符号 64 位输入；0 不在定义域内。

本地双版本验证包含 $1\ldots20000$ 的线性筛交叉检查、小范围直接计数互质整数，
以及重复质因子、大质数平方、$2^{64}-1$ 和已独立认证的大素数。
普通与 ASan/UBSan 测试均通过；P4718 的 AC 只验证底层分解用途，不是此欧拉函数的新 AC。
来源：kuangbin 2.7.1、2.7.3；WIDA 的单个数欧拉函数。

\section{线性筛：三张表的约定}\label{knowledge-linear-sieve}\index{linear sieve}
\texttt{LinearSieve} 同时构造最小质因子、欧拉函数和莫比乌斯函数。
内层遇到第一个整除 $i$ 的素数后停止，因此每个合数 $j$ 只由
$p=\operatorname{lp}(j)$ 与 $i=j/p$ 生成一次，时间为 $O(n)$。
若 $p\mid i$，则 $\varphi(ip)=p\varphi(i)$、$\mu(ip)=0$；
否则二者分别为 $(p-1)\varphi(i)$ 和 $-\mu(i)$。

判素使用 $x\ge2$ 且 $\operatorname{lp}(x)=x$，不能漏掉前半条件。
分解表内正整数时，反复取 $p=\operatorname{lp}(x)$、令 $x\gets x/p$，直到 $x=1$；
重复因子会重复输出。0 不可分解，1 的质因子列表为空。
上限要求 $0\le n<\texttt{INT\_MAX}$，避免构造时 $n+1$ 溢出，并保证内存可用。
筛乘积前先检查 $p\le n/i$，不直接计算可能溢出的 $ip$。

双版本验证至 $n=10^7$，包括小前缀、试除最小质因子、因子还原和逐项约数和证书：
\[
\sum_{d\mid x}\varphi(d)=x,\qquad\sum_{d\mid x}\mu(d)=[x=1].
\]
这些恒等式在 $d=x$ 项系数为 1，可按 $x$ 递增唯一确定两张表。
普通与 ASan/UBSan 均通过；不据此宣称莫比乌斯反演的所有应用已验证。

来源：kuangbin 2.1.1、2.1.2、2.7.2、2.7.4、2.13.1；WIDA 的全部欧拉函数表。

WIDA 的“使用莫比乌斯反演求解欧拉函数”先存储全部约数，再使用
\[
\varphi(n)=n-\sum_{d\mid n,\ d<n}\varphi(d).
\]
这与 $\varphi*1=\operatorname{id}$ 等价；再卷积 $\mu$ 可得
$\varphi(n)=\sum_{d\mid n}\mu(d)(n/d)$。
本库用线性筛提供同一张表，无需额外保存约数列表。
百万范围另分别运行直接卷积和真约数扣减，逐项核对两套筛表；两种方法都不复用线性筛的递推分支。


\section{区间筛：只保存目标区间}\label{knowledge-segmented-sieve}\index{segmented sieve}\index{区间筛素数}
实现见第~\pageref{compact-segmented_primes}~页。
设 $0\le L\le R\le\texttt{INT\_MAX}$，先筛出不超过 $\lfloor\sqrt R\rfloor$ 的素数，
再对每个素数 $p$ 从
\[
\max\bigl(p^2,\lceil L/p\rceil p\bigr)
\]
开始标记区间内的倍数。任何合数都有不超过自身平方根的素因子，所以不会遗漏；
从 $p^2$ 开始也不会误删 $p$ 自身。输出跳过 0 和 1，其余未标记数恰为素数。
底筛与区间标记共用同一个外层素数循环，但使用两张独立的表。

令 $W=R-L+1$，空间为 $O(\sqrt R+W)$，包含返回的素数数组；
时间为 $O((\sqrt R+W)\log\log(R+3))$。函数每次调用都重建底筛，适合题目中的少量区间查询。
端点虽然是 \texttt{int}，向上取整的加法、$p^2$、标记变量和输出遍历变量均用
\texttt{long long}；否则 $R=\texttt{INT\_MAX}$ 时最后一次递增也会溢出。
区间长度仍受可用内存约束，不适合一次保存整个 32 位值域。

POJ 2689 与 UVa 10140 的驱动读到 EOF，扫描相邻素数差。
最近、最远都只在严格更优时更新，距离并列保留最先出现的一对；少于两个素数时输出题目指定句子。
两套实现以试除和 Miller--Rabin 判素进行独立对拍，并检查平方边界、单点区间及
右端点为 \texttt{INT\_MAX} 的百万长度区间；完整驱动验证 208 组输入，含官方样例。
普通测试与 ASan/UBSan 均通过；尚无本库在线 AC。

来源：kuangbin 2.1.3；\url{https://onlinejudge.org/external/101/10140.pdf}。

